\chapter{HIỆN THỰC HÓA VÀ TRIỂN KHAI}

\section{Cấu trúc mã nguồn dự án}

Dự án được tổ chức theo mô hình monorepo với ba thư mục chính: \texttt{backend/} chứa Backend Orchestrator, \texttt{frontend/} chứa Frontend Dashboard, và \texttt{infrastructure/} chứa cấu hình Docker Compose và Monitoring Stack.

\begin{lstlisting}[caption={Cấu trúc thư mục dự án}, label={lst:project_structure}, language={}]
Layer1/
  backend/
    src/
      auth/            # JWT + Google OAuth
      subnets/         # Subnet CRUD + Lifecycle
      validators/      # Docker container mgt
      node-status/     # Status polling + WS Gateway
      rpc-proxy/       # JSON-RPC forwarding
      contracts/       # Smart contract deployment
      benchmarks/      # Perf testing
      pchain/          # P-Chain API queries
      prisma/          # Database access layer
      network-runner/  # Avalanche Network Runner
    prisma/
      schema.prisma    # DB schema definition
  frontend/
    src/
      app/             # Next.js App Router
        login/         # Login page
        register/      # Registration page
      components/v2/   # UI components (TopBar, etc.)
      context/         # AuthContext, NetworkContext
      lib/             # API client, helpers
  infrastructure/
    docker-compose.yml # Postgres + Prometheus + Grafana
    prometheus.yml     # Scrape config
\end{lstlisting}

\section{Hiện thực Backend Orchestrator}

\subsection{Cấu hình và khởi tạo ứng dụng}

File \texttt{main.ts} là điểm khởi đầu của toàn bộ Backend:

\begin{lstlisting}[caption={main.ts --- Khởi tạo ứng dụng NestJS}, label={lst:main_ts}, language=Java]
async function bootstrap() {
  const app = await NestFactory.create(AppModule, {
    logger: ['log', 'warn', 'error'],
    // Loai bo DEBUG de giam noise khi node offline
  });
  app.enableCors({
    origin: process.env.FRONTEND_URL || '*',
    credentials: true,
  });
  await app.listen(process.env.PORT || 4000);
  Logger.log(`API running on http://localhost:4000`,
    'Bootstrap');
}
\end{lstlisting}

Em cấu hình logger ở mức \texttt{['log', 'warn', 'error']} thay vì mặc định (bao gồm cả debug) để giữ console sạch khi Avalanche node chưa khởi động. Kỹ thuật này gọi là ``Graceful Degradation'' --- hệ thống vẫn khởi động và phục vụ được các chức năng không phụ thuộc node (đăng ký, đăng nhập, quản lý Subnet trong database) ngay cả khi node offline.

\subsection{Module xác thực (AuthModule)}

Module xác thực là thành phần then chốt đảm bảo an toàn cho toàn bộ hệ thống. Em triển khai hai phương thức đăng nhập song song: Email/Password (Local Strategy) và Google OAuth 2.0.

\subsubsection{Cấu trúc AuthModule}

\begin{lstlisting}[caption={AuthModule --- Đăng ký các thành phần xác thực}, label={lst:auth_module}, language=Java]
@Module({
  imports: [
    PrismaModule,
    JwtModule.register({
      secret: process.env.JWT_SECRET,
      signOptions: { expiresIn: '24h' },
    }),
    PassportModule,
  ],
  controllers: [AuthController],
  providers: [
    AuthService,
    JwtStrategy,     // Xac thuc JWT token
    GoogleStrategy,  // Xac thuc Google OAuth
  ],
  exports: [AuthService],
})
export class AuthModule {}
\end{lstlisting}

NestJS Dependency Injection Container tự động quản lý vòng đời của các Provider. Khi \texttt{SubnetsController} cần kiểm tra JWT, NestJS inject \texttt{JwtStrategy} (đã đăng ký trong AuthModule) vào pipeline mà không cần bất kỳ cấu hình thủ công nào.

\subsubsection{Đăng ký và đăng nhập bằng Email}

Quy trình đăng ký mã hóa mật khẩu bằng \texttt{bcrypt} với salt round = 10 trước khi lưu:

\begin{lstlisting}[caption={AuthService --- Đăng ký và đăng nhập}, label={lst:auth_register}, language=Java]
async register(dto: RegisterDto) {
  const existing = await this.prisma.user.findUnique({
    where: { email: dto.email },
  });
  if (existing) {
    throw new ConflictException('Email already in use');
  }
  const hashedPassword = await bcrypt.hash(
    dto.password, 10
  );
  const user = await this.prisma.user.create({
    data: {
      email: dto.email,
      password: hashedPassword,
      name: dto.name,
    },
  });
  const { password, ...result } = user;
  return result;
}

async login(dto: LoginDto) {
  const user = await this.prisma.user.findUnique({
    where: { email: dto.email },
  });
  if (!user || !await bcrypt.compare(
    dto.password, user.password
  )) {
    throw new UnauthorizedException('Invalid credentials');
  }
  const payload = { sub: user.id, email: user.email };
  return {
    access_token: this.jwtService.sign(payload),
    user: { id: user.id, email: user.email,
            name: user.name },
  };
}
\end{lstlisting}

Lưu ý về bảo mật: hàm \texttt{bcrypt.compare()} sử dụng constant-time comparison để tránh timing attack --- kẻ tấn công không thể đoán mật khẩu bằng cách đo thời gian phản hồi.

\subsubsection{Tích hợp Google OAuth 2.0}

\begin{lstlisting}[caption={GoogleStrategy --- Xác thực và tự động tạo user}, label={lst:google_strategy}, language=Java]
@Injectable()
export class GoogleStrategy extends PassportStrategy(
  Strategy, 'google'
) {
  constructor(private readonly prisma: PrismaService) {
    super({
      clientID: process.env.GOOGLE_CLIENT_ID,
      clientSecret: process.env.GOOGLE_CLIENT_SECRET,
      callbackURL: process.env.GOOGLE_CALLBACK_URL,
      scope: ['email', 'profile'],
    });
  }

  async validate(accessToken, refreshToken,
    profile, done
  ) {
    const email = profile.emails?.[0]?.value;
    let user = await this.prisma.user.findUnique({
      where: { email },
    });
    if (!user) {
      // Auto-create account from Google profile
      user = await this.prisma.user.create({
        data: {
          email,
          name: profile.displayName,
          password: await bcrypt.hash(
            Math.random().toString(36), 10
          ), // Random password (Google users login via OAuth)
        },
      });
    }
    const { password, ...result } = user;
    done(null, result);
  }
}
\end{lstlisting}

Callback handler redirect về Frontend kèm JWT:

\begin{lstlisting}[caption={AuthController --- Google OAuth callback}, label={lst:google_callback}, language=Java]
@Get('google/callback')
@UseGuards(AuthGuard('google'))
async googleCallback(@Req() req, @Res() res) {
  const jwt = await this.authService
    .issueTokenForOAuthUser(req.user);
  const userJson = encodeURIComponent(
    JSON.stringify(req.user)
  );
  res.redirect(
    `${process.env.FRONTEND_URL}/login` +
    `?token=${jwt.access_token}&user=${userJson}`
  );
}
\end{lstlisting}

\subsubsection{Bảo vệ endpoint bằng JwtAuthGuard}

Mọi endpoint nhạy cảm đều được bảo vệ bởi decorator \texttt{@UseGuards(AuthGuard('jwt'))}. JwtStrategy giải mã token, trích xuất user ID, và inject vào \texttt{req.user}:

\begin{lstlisting}[caption={JwtStrategy --- Xác minh JWT token}, label={lst:jwt_strategy}, language=Java]
@Injectable()
export class JwtStrategy extends PassportStrategy(Strategy) {
  constructor(private prisma: PrismaService) {
    super({
      jwtFromRequest: ExtractJwt.fromAuthHeaderAsBearerToken(),
      secretOrKey: process.env.JWT_SECRET,
    });
  }
  async validate(payload: { sub: string; email: string }) {
    const user = await this.prisma.user.findUnique({
      where: { id: payload.sub },
    });
    if (!user) throw new UnauthorizedException();
    return { id: user.id, email: user.email,
             name: user.name };
  }
}
\end{lstlisting}


\subsection{Module quản lý Subnet (SubnetsModule)}

\subsubsection{Vòng đời Subnet}

Mỗi Subnet trong hệ thống trải qua các trạng thái: \texttt{creating} $\rightarrow$ \texttt{running} $\rightarrow$ \texttt{stopped} $\rightarrow$ \texttt{deleted}. SubnetsService quản lý toàn bộ vòng đời này:

\begin{figure}[H]
\centering
\fbox{\parbox{0.85\textwidth}{\centering\vspace{2cm}\textit{[Hình 4.1: State Machine Diagram --- Vòng đời trạng thái Subnet]}\vspace{2cm}}}
\caption{Biểu đồ trạng thái vòng đời Subnet}
\label{fig:subnet_lifecycle}
\end{figure}

\subsubsection{Giải thuật cấp phát Port tự động}

Khi tạo Docker container cho validator, hệ thống cần cấp phát cổng mạng tránh xung đột:

\begin{lstlisting}[caption={Thuật toán cấp phát Port tự động}, label={lst:port_allocation}, language=Java]
private async findAvailablePort(
  startPort: number = 19650
): Promise<number> {
  const usedPorts = new Set<number>();

  // Thu thap tat ca port dang su dung tu Docker
  const containers = await this.docker.listContainers();
  for (const container of containers) {
    for (const port of container.Ports || []) {
      if (port.PublicPort) {
        usedPorts.add(port.PublicPort);
      }
    }
  }

  // Thu thap port tu database
  const validators = await this.prisma.validator
    .findMany({ select: { port: true } });
  for (const v of validators) {
    if (v.port) usedPorts.add(v.port);
  }

  // Tim port ranh trong dai [startPort, startPort+1000]
  for (let port = startPort;
       port < startPort + 1000; port++) {
    if (!usedPorts.has(port)) {
      return port;
    }
  }
  throw new Error('No available ports in range');
}
\end{lstlisting}

Thuật toán có độ phức tạp $O(N + P)$ với $N$ là số container đang chạy và $P$ là kích thước dải port. Trong thực tế, $N < 50$ nên thuật toán chạy gần tức thì. Thuật toán kiểm tra cả Docker lẫn database để xử lý trường hợp container đã dừng nhưng port vẫn bị giữ trong database.

\subsubsection{Kiểm tra quyền sở hữu}

Mỗi thao tác thay đổi Subnet đều phải kiểm tra quyền sở hữu (ownership check pattern):

\begin{lstlisting}[caption={Pattern kiểm tra quyền sở hữu Subnet}, label={lst:ownership}, language=Java]
private async validateOwnership(
  subnetId: string, userId: string
): Promise<Network> {
  const network = await this.prisma.network.findUnique({
    where: { id: subnetId },
    include: { validators: true },
  });
  if (!network) {
    throw new NotFoundException('Subnet not found');
  }
  if (network.userId !== userId) {
    throw new ForbiddenException(
      'You do not own this subnet'
    );
  }
  return network;
}

// Su dung trong moi endpoint:
async stop(subnetId: string, userId: string) {
  const network = await this.validateOwnership(
    subnetId, userId
  );
  // Dung tat ca validator container...
  for (const validator of network.validators) {
    if (validator.containerId) {
      await this.docker.getContainer(
        validator.containerId
      ).stop();
    }
  }
  await this.prisma.network.update({
    where: { id: subnetId },
    data: { status: 'stopped' },
  });
}
\end{lstlisting}


\subsection{Module giám sát (NodeStatusModule)}

\subsubsection{Kiểm tra khả dụng trước khi truy vấn}

Vấn đề thực tế em gặp phải: khi Avalanche node offline, thư viện \texttt{ethers.js} sẽ tự động retry tạo \texttt{JsonRpcProvider}, in ra hàng dòng lỗi ``JsonRpcProvider failed to detect network and cannot start up; retry in 1s'' liên tục. Để giải quyết, em triển khai cơ chế pre-check bằng một HTTP probe nhẹ:

\begin{lstlisting}[caption={isRpcReachable --- Kiểm tra khả dụng trước khi tạo Provider}, label={lst:rpc_precheck}, language=Java]
private async isRpcReachable(
  rpcUrl: string
): Promise<boolean> {
  try {
    const controller = new AbortController();
    const timeout = setTimeout(
      () => controller.abort(), 2000
    );
    const res = await fetch(rpcUrl, {
      method: 'POST',
      headers: { 'Content-Type': 'application/json' },
      body: JSON.stringify({
        jsonrpc: '2.0', id: 1,
        method: 'eth_blockNumber', params: []
      }),
      signal: controller.signal,
    });
    clearTimeout(timeout);
    return res.ok;
  } catch {
    return false;
  }
}
\end{lstlisting}

Khi \texttt{isRpcReachable} trả về \texttt{false}, hệ thống bỏ qua mọi truy vấn ethers.js và trả về trạng thái offline yên lặng --- không có bất kỳ lỗi nào trong log. Kỹ thuật này áp dụng cho cả ba phương thức: \texttt{fetchStatus()}, \texttt{fetchDashboardStats()}, và \texttt{fetchRecentBlocks()}.

\subsubsection{WebSocket Gateway --- Phát sóng real-time}

NodeStatusGateway sử dụng Socket.IO để duy trì kết nối hai chiều với Frontend:

\begin{lstlisting}[caption={NodeStatusGateway --- Quản lý kết nối và phát sóng}, label={lst:ws_gateway}, language=Java]
@WebSocketGateway({
  cors: { origin: '*' },
  namespace: 'node-status'
})
export class NodeStatusGateway
  implements OnGatewayConnection, OnGatewayDisconnect
{
  private globalInterval: NodeJS.Timeout | null = null;
  private clientIntervals: Map<string, NodeJS.Timeout>
    = new Map();

  handleConnection(client: Socket) {
    if (!this.globalInterval) {
      this.startBroadcasting();
    }
  }

  handleDisconnect(client: Socket) {
    // Cleanup interval cua client nay
    const interval = this.clientIntervals.get(client.id);
    if (interval) {
      clearInterval(interval);
      this.clientIntervals.delete(client.id);
    }
    // Neu khong con client nao, dung global interval
    if (this.server.sockets.sockets.size === 0) {
      clearInterval(this.globalInterval);
      this.globalInterval = null;
    }
  }

  private startBroadcasting() {
    this.globalInterval = setInterval(async () => {
      const status = await this.nodeStatusService
        .fetchStatus();
      this.server.emit('status', status);
    }, 5000);
  }

  @SubscribeMessage('subscribe_chain')
  handleChainSubscribe(client: Socket, rpcUrl: string) {
    // Tao interval rieng cho client nay
    const interval = setInterval(async () => {
      const stats = await this.nodeStatusService
        .fetchChainStats(rpcUrl);
      client.emit('chain-stats', stats);
    }, 3000);
    this.clientIntervals.set(client.id, interval);
  }
}
\end{lstlisting}

Thiết kế này có hai điểm kỹ thuật quan trọng. Thứ nhất, \texttt{globalInterval} chỉ chạy khi có ít nhất một client kết nối, tránh lãng phí CPU khi không ai xem Dashboard. Thứ hai, \texttt{clientIntervals} Map quản lý interval riêng cho từng client subscription, và tự động cleanup khi client ngắt kết nối để tránh memory leak.


\subsection{Module RPC Proxy}

RpcProxyService đóng vai trò proxy trung gian giữa Frontend và Avalanche node. Khi Frontend gọi một RPC request (ví dụ \texttt{eth\_getBalance}), request được chuyển tiếp qua Backend thay vì gọi trực tiếp node. Cách tiếp cận này mang lại:

\begin{itemize}
    \item \textbf{Bảo mật:} Node RPC endpoint không bị expose trực tiếp ra internet.
    \item \textbf{Xử lý lỗi tập trung:} Backend bắt các lỗi ECONNREFUSED hoặc 503 khi node offline, trả về JSON-RPC error response chuẩn thay vì để Frontend bị crash.
    \item \textbf{Rate limiting:} Có thể áp đặt giới hạn số request/giây cho mỗi user.
\end{itemize}

\begin{lstlisting}[caption={RpcProxyService --- Chuyển tiếp JSON-RPC}, label={lst:rpc_proxy}, language=Java]
@Post('proxy')
async proxyRpc(
  @Body() body: any,
  @Query('rpcUrl') rpcUrl: string
) {
  try {
    const response = await axios.post(rpcUrl, body, {
      timeout: 10000,
      headers: { 'Content-Type': 'application/json' },
    });
    return response.data;
  } catch (error) {
    if (error.code === 'ECONNREFUSED'
      || error.response?.status === 503) {
      // Node offline - tra ve loi RPC chuan
      return {
        jsonrpc: '2.0', id: body.id || 1,
        error: {
          code: -32603,
          message: 'Node is offline'
        }
      };
    }
    throw error;
  }
}
\end{lstlisting}

\subsection{Module quản lý Validator (ValidatorsModule)}

ValidatorsService quản lý các Docker container chứa AvalancheGo node. Khi hệ thống khởi động, Service thực hiện quy trình \texttt{recoverDockerPorts()} để đồng bộ lại trạng thái container với database:

\begin{lstlisting}[caption={ValidatorsService --- Phục hồi trạng thái khi khởi động}, label={lst:validator_recover}, language=Java]
async recoverDockerPorts() {
  // Kiem tra Docker co dang chay khong
  try {
    await execPromise('docker info');
  } catch {
    this.logger.debug(
      'Docker not available, skipping recovery'
    );
    return;
  }

  const validators = await this.prisma.validator
    .findMany({
      where: { containerId: { not: null } },
    });

  for (const val of validators) {
    try {
      const output = await execPromise(
        `docker port ${val.containerId} 9650`
      );
      const port = parseInt(
        output.split(':').pop().trim()
      );
      await this.prisma.validator.update({
        where: { id: val.id },
        data: { port },
      });
    } catch (err) {
      if (err.message.includes('No such container')) {
        // Container da bi xoa, xoa containerId
        await this.prisma.validator.update({
          where: { id: val.id },
          data: { containerId: null },
        });
      }
    }
  }
}
\end{lstlisting}

Đoạn code trên giải quyết vấn đề thực tế: khi hệ thống restart, Docker container có thể vẫn đang chạy nhưng port mapping trong database đã cũ. Quy trình recovery đồng bộ lại port thực tế, và tự động dọn dẹp các bản ghi trỏ tới container đã không còn tồn tại.


\section{Hiện thực Frontend Dashboard}

\subsection{Quản lý trạng thái xác thực (AuthContext)}

Frontend sử dụng React Context API để quản lý trạng thái xác thực xuyên suốt ứng dụng. Thiết kế AuthContext tuân theo nguyên tắc Single Source of Truth:

\begin{lstlisting}[caption={AuthContext --- Quản lý trạng thái xác thực}, label={lst:auth_context}, language=Java]
export function AuthProvider({ children }) {
  const [user, setUser] = useState(null);
  const [token, setToken] = useState(null);
  const [loading, setLoading] = useState(true);

  useEffect(() => {
    // Khoi phuc phien dang nhap tu localStorage
    const savedToken = localStorage.getItem('jwt_token');
    const savedUser = localStorage.getItem('user_data');
    if (savedToken && savedUser) {
      setToken(savedToken);
      setUser(JSON.parse(savedUser));
    }
    setLoading(false); // Ket thuc kiem tra
  }, []);

  const login = (newToken, userData) => {
    setToken(newToken);
    setUser(userData);
    localStorage.setItem('jwt_token', newToken);
    localStorage.setItem('user_data',
      JSON.stringify(userData));
  };

  const logout = () => {
    setToken(null);
    setUser(null);
    localStorage.removeItem('jwt_token');
    localStorage.removeItem('user_data');
  };

  return (
    <AuthContext.Provider value={{
      user, token, loading, login, logout
    }}>
      {children}
    </AuthContext.Provider>
  );
}
\end{lstlisting}

Trạng thái \texttt{loading} rất quan trọng: nó ngăn chặn hiện tượng ``flash'' --- trang Dashboard hiện ra rồi ngay lập tức redirect về Login trong khi Context đang đọc localStorage. Bằng cách hiển thị loading spinner cho đến khi kiểm tra xong, trải nghiệm người dùng trở nên mượt mà hơn.

\subsection{Tích hợp JWT vào API Client}

Mọi yêu cầu API từ Frontend đều tự động đính kèm JWT token qua wrapper function \texttt{fetchJson}:

\begin{lstlisting}[caption={API Client --- Tự động đính kèm JWT}, label={lst:api_client}, language=Java]
const API_BASE = process.env.NEXT_PUBLIC_API_URL
  || 'http://localhost:4000';

async function fetchJson(path, options = {}) {
  const token = localStorage.getItem('jwt_token');
  const headers = {
    'Content-Type': 'application/json',
    ...options.headers,
  };
  if (token) {
    headers['Authorization'] = `Bearer ${token}`;
  }
  const url = `${API_BASE}${path}`;
  const response = await fetch(url, {
    ...options, headers
  });
  if (response.status === 401) {
    // Token het han, tu dong logout
    localStorage.removeItem('jwt_token');
    window.location.href = '/login';
    return;
  }
  if (!response.ok) {
    const err = await response.json().catch(() => ({}));
    throw new Error(
      err.message || `HTTP ${response.status}`
    );
  }
  return response.json();
}
\end{lstlisting}

Điểm đáng chú ý: khi nhận được HTTP 401 (Unauthorized), API client tự động xóa token cũ và redirect về trang Login thay vì hiển thị lỗi khó hiểu cho người dùng.

\subsection{Xử lý Google OAuth Callback}

Trang Login sử dụng \texttt{useEffect} kết hợp \texttt{useSearchParams} để xử lý callback từ Google:

\begin{lstlisting}[caption={Xử lý Google OAuth callback trên Frontend}, label={lst:oauth_callback}, language=Java]
function LoginCallbackHandler() {
  const searchParams = useSearchParams();
  const { login } = useAuth();
  const router = useRouter();

  useEffect(() => {
    const token = searchParams.get('token');
    const userParam = searchParams.get('user');
    if (token && userParam) {
      try {
        const user = JSON.parse(
          decodeURIComponent(userParam)
        );
        login(token, user);
        router.push('/');
      } catch (e) {
        console.error('OAuth callback failed:', e);
      }
    }
  }, [searchParams]);

  return null; // Renderer invisible
}

// Wrap voi Suspense vi useSearchParams la client-side
export default function LoginPage() {
  return (
    <Suspense fallback={<Loading />}>
      <LoginCallbackHandler />
      <LoginForm />
    </Suspense>
  );
}
\end{lstlisting}

Component \texttt{LoginCallbackHandler} phải được wrap trong \texttt{Suspense} vì Next.js App Router yêu cầu \texttt{useSearchParams} chỉ được sử dụng trong Client Component boundary. Đây là một chi tiết kỹ thuật mà em đã trải qua quá trình debug thực tế mới phát hiện.

\subsection{Giao diện người dùng}

\begin{figure}[H]
\centering
\fbox{\parbox{0.85\textwidth}{\centering\vspace{3cm}\textit{[Hình 4.2: Giao diện trang đăng nhập --- Google OAuth + Email/Password]}\vspace{3cm}}}
\caption{Giao diện trang đăng nhập}
\label{fig:login_page}
\end{figure}

\begin{figure}[H]
\centering
\fbox{\parbox{0.85\textwidth}{\centering\vspace{3cm}\textit{[Hình 4.3: Dashboard chính --- Hiển thị trạng thái node, block height, TPS, peer count]}\vspace{3cm}}}
\caption{Giao diện Dashboard chính}
\label{fig:dashboard}
\end{figure}

\begin{figure}[H]
\centering
\fbox{\parbox{0.85\textwidth}{\centering\vspace{3cm}\textit{[Hình 4.4: Wizard tạo Subnet --- Bước 1: Nhập tên và chọn VM type]}\vspace{3cm}}}
\caption{Wizard tạo Subnet --- Bước 1}
\label{fig:wizard_step1}
\end{figure}

\begin{figure}[H]
\centering
\fbox{\parbox{0.85\textwidth}{\centering\vspace{3cm}\textit{[Hình 4.5: Giao diện danh sách Subnet với trạng thái và thao tác]}\vspace{3cm}}}
\caption{Trang quản lý Subnet}
\label{fig:subnets_page}
\end{figure}


\section{Cấu hình hạ tầng và triển khai}

\subsection{Docker Compose --- Dịch vụ hỗ trợ}

Các dịch vụ nền (PostgreSQL, Prometheus, Grafana) được quản lý bằng Docker Compose:

\begin{lstlisting}[caption={docker-compose.yml --- Dịch vụ hạ tầng}, label={lst:docker_compose}, language={}]
version: '3.8'
services:
  postgres:
    image: postgres:15-alpine
    environment:
      POSTGRES_USER: layer1
      POSTGRES_PASSWORD: ${DB_PASSWORD}
      POSTGRES_DB: layer1_db
    ports:
      - "5432:5432"
    volumes:
      - pgdata:/var/lib/postgresql/data

  prometheus:
    image: prom/prometheus:latest
    ports:
      - "9090:9090"
    volumes:
      - ./prometheus.yml:/etc/prometheus/prometheus.yml
    command:
      - '--config.file=/etc/prometheus/prometheus.yml'
      - '--storage.tsdb.retention.time=15d'

  grafana:
    image: grafana/grafana:latest
    ports:
      - "3001:3000"
    depends_on:
      - prometheus
    environment:
      GF_SECURITY_ADMIN_PASSWORD: ${GRAFANA_PASSWORD}

volumes:
  pgdata:
\end{lstlisting}

\subsection{Prometheus --- Cấu hình thu thập metrics}

\begin{lstlisting}[caption={prometheus.yml --- Cấu hình scrape targets}, label={lst:prometheus_config}, language={}]
global:
  scrape_interval: 15s
  evaluation_interval: 15s

scrape_configs:
  - job_name: 'avalanche-primary'
    static_configs:
      - targets: ['host.docker.internal:9650']
    metrics_path: '/ext/metrics'

  - job_name: 'avalanche-subnets'
    file_sd_configs:
      - files:
          - '/etc/prometheus/subnet_targets.json'
        refresh_interval: 30s
\end{lstlisting}

Cấu hình \texttt{file\_sd\_configs} cho phép Prometheus tự động phát hiện các Subnet mới khi Backend cập nhật file \texttt{subnet\_targets.json}. Đây là kỹ thuật Service Discovery nâng cao, cho phép monitoring stack tự động mở rộng theo số lượng Subnet mà không cần restart Prometheus.
